% !TeX root = ../main.tex
% !TeX program = pdflatex
% !TeX encoding = UTF-8

\chapter{Hạn chế và hướng phát triển}
\label{chap:limitations}

\section{Khó khăn trong quá trình phát triển}

Khó khăn chính là thống nhất dữ liệu giữa công cụ tạo mạng đường, backend và
frontend. ID của đỉnh và cạnh phải ổn định để trace thuật toán khớp chính xác
với hình vẽ. Nhóm cũng phải thiết kế một định dạng kết quả chung cho nhiều
thuật toán có cách quản lý frontier khác nhau, đồng thời bảo đảm giao diện phát
lại được toàn bộ quá trình tìm kiếm.

Một thách thức khác là xây dựng chi phí có ý nghĩa khi khoảng cách, thời gian,
ùn tắc và rủi ro có thang đo khác nhau. Việc chuẩn hóa và chọn bốn bộ trọng số
giúp người dùng hiểu tiêu chí tối ưu, nhưng vẫn cần kiểm thử nhiều trường hợp để
tránh kết quả trái trực giác. Chức năng đa điểm còn phải cân bằng giữa nghiệm
chính xác của Brute Force và tốc độ của Nearest Neighbor.

\section{Hạn chế của dữ liệu, hàm chi phí và thuật toán}

\begin{itemize}[itemsep=2pt]
  \item \textbf{Tập dữ liệu:} chỉ có 25 địa điểm và 114 cạnh giữa các điểm quan
  tâm, không phải toàn bộ mạng giao lộ Đà Lạt. Hình học hiển thị cạnh được đơn
  giản thành đoạn thẳng.
  \item \textbf{Điều kiện giao thông:} năm kịch bản được sinh mô phỏng và cố
  định trong mỗi lần chạy, không cập nhật theo thời gian thực. Đóng đường, mưa,
  sương mù và thi công là giá trị giả lập.
  \item \textbf{Hàm chi phí:} chuẩn hóa min--max phụ thuộc tập cạnh còn mở của
  từng kịch bản. Các trọng số được định trước và chưa cho phép người dùng tự
  điều chỉnh; chi phí cũng chưa phản ánh độ dốc, phí đường hoặc loại phương tiện.
  \item \textbf{Heuristic:} khi tối ưu \texttt{route\_cost} hoặc rủi ro, A* dùng
  heuristic 0 để bảo đảm không đánh giá vượt, vì vậy không mở rộng ít đỉnh hơn UCS.
  \item \textbf{Đa điểm:} Nearest Neighbor không bảo đảm tối ưu toàn cục; Brute
  Force cho nghiệm chính xác nhưng tăng theo giai thừa và bị giới hạn tám đích.
  \item \textbf{Đánh giá hiệu năng:} dataset nhỏ và bảng thời gian mới dựa trên
  một lần chạy, nên chưa đủ để kết luận về tốc độ trên mạng đường quy mô lớn.
\end{itemize}

\section{Hướng phát triển}

\begin{itemize}[itemsep=2pt]
  \item Tích hợp API bản đồ và giao thông thời gian thực để cập nhật tốc độ,
  ùn tắc, đóng đường và hình học tuyến.
  \item Mở rộng mạng đến cấp giao lộ, bổ sung độ dốc và xây dựng profile riêng
  cho ô tô, xe máy, xe đạp và đi bộ.
  \item Cho phép người dùng chỉnh trọng số, lưu cấu hình và so sánh nhiều tuyến
  ứng viên trên cùng một màn hình.
  \item Dùng các phương pháp như 2-opt, genetic algorithm hoặc ant colony để
  cải thiện bài toán đa điểm lớn; xa hơn là mở rộng sang nhiều phương tiện.
  \item Thực hiện benchmark lặp lại trên nhiều cặp nguồn--đích và nhiều quy mô
  đồ thị, báo cáo trung bình, độ lệch chuẩn và chất lượng nghiệm.
  \item Triển khai backend/frontend trên hạ tầng công khai và bổ sung cache,
  giám sát lỗi cùng kiểm thử tích hợp đầu cuối.
\end{itemize}
